\section{Resumen}

Este proyecto desarrolla desde cero algoritmos genéticos para optimizar redes neuronales en dos vertientes, comparándolos siempre con líneas base en igualdad de presupuesto. En la primera, se evolucionan directamente los pesos de una red neuronal de control en CartPole-v1, entorno cuya función de recompensa no admite gradientes útiles. El algoritmo (selección por torneo y ruleta, elitismo, cruces de un punto, uniforme y BLX-$\alpha$, mutación gaussiana y mixta) resuelve la tarea en 5 de 5 semillas con 470 evaluaciones de media, frente a 700 de la búsqueda aleatoria y 3933 episodios de REINFORCE, que solo alcanza el objetivo en 3 de 5 ejecuciones. La comparación de operadores sitúa la ruleta (220 evaluaciones) y BLX-$\alpha$ (310) por delante del cruce de un punto (470) y del uniforme (550). En la segunda vertiente, el algoritmo actúa como optimizador de caja negra de la arquitectura y los hiperparámetros de una red entrenada con descenso por gradiente sobre el conjunto tabular UCI Letter Recognition, con un cromosoma mixto de 10 genes: alcanza 0,9619 de exactitud en validación frente a 0,9501 de la búsqueda aleatoria y 0,8733 de la configuración por defecto, con 0,9612 en test. La evaluación paralela de poblaciones, verificada con semillas preasignadas, logra una aceleración de 4,76$\times$ en 12 núcleos. Los resultados concluyen que la evolución es más eficiente y fiable que la suerte y que el gradiente en este tipo de tareas, con las salvedades estadísticas de 5 semillas y de un solo problema por vertiente documentadas en la discusión.
